\chapter{KINEMATICS 运动学}
\intro{盖周天之变，先天领周天，化吾为王。}
\section{简单转动(Simple Rotation)}
$\vec{a}$是坐标系$\mathbb{A}$中的一个向量。$\vec{b}$是坐标系$\mathbb{B}$中的一个向量。
且坐标系A和B之间存在一个旋转关系。则有：
\begin{myformula}[title={Rodrigues 旋转公式}]
    \vec{b}=\vec{a}\cos\theta-\vec{a}\times\vec{\lambda}\sin\theta+\vec{\lambda}(\vec{\lambda}\cdot\vec{a})(1-\cos\theta)
\end{myformula}
其中：
\begin{itemize}
    \item $\theta$ 是旋转角度
    \item $\vec{\lambda}$ 是旋转轴的单位向量
\end{itemize}

\begin{proof}[Rodrigues 旋转公式的推导]
在坐标系 $\mathbb{A}$ 中，以转轴 $\vec{\lambda}$ 为第一个基向量，构造右手标准正交基 $\{\vec{\lambda},\,\vec{\mu}_A,\,\vec{\nu}_A\}$，其中 $\vec{\nu}_A = \vec{\lambda}\times\vec{\mu}_A$。将 $\vec{a}$ 在此基下展开：
\[
\vec{a} = a_\lambda\vec{\lambda} + a_\mu\vec{\mu}_A + a_\nu\vec{\nu}_A,
\qquad
a_\lambda = \vec{a}\cdot\vec{\lambda},\;
a_\mu = \vec{a}\cdot\vec{\mu}_A,\;
a_\nu = \vec{a}\cdot\vec{\nu}_A
\]

坐标系 $\mathbb{B}$ 由 $\mathbb{A}$ 绕 $\vec{\lambda}$ 旋转角度 $\theta$ 得到，其基向量 $\{\vec{\lambda},\,\vec{\mu}_B,\,\vec{\nu}_B\}$ 在 $\mathbb{A}$ 系中表为：
\begin{align*}
    \vec{\mu}_B &= \vec{\mu}_A\cos\theta + \vec{\nu}_A\sin\theta \\
    \vec{\nu}_B &= -\vec{\mu}_A\sin\theta + \vec{\nu}_A\cos\theta
\end{align*}

同一物理向量在 $\mathbb{B}$ 系中的坐标与 $\mathbb{A}$ 系相同（均为 $a_\lambda,a_\mu,a_\nu$），故：
\begin{align*}
    \vec{b} &= a_\lambda\vec{\lambda} + a_\mu\vec{\mu}_B + a_\nu\vec{\nu}_B \\
            &= a_\lambda\vec{\lambda} + a_\mu(\vec{\mu}_A\cos\theta + \vec{\nu}_A\sin\theta)
                                     + a_\nu(-\vec{\mu}_A\sin\theta + \vec{\nu}_A\cos\theta) \\
            &= a_\lambda\vec{\lambda} + \cos\theta\,(a_\mu\vec{\mu}_A + a_\nu\vec{\nu}_A)
                                     + \sin\theta\,(a_\mu\vec{\nu}_A - a_\nu\vec{\mu}_A)
\end{align*}

由 $\vec{\nu}_A = \vec{\lambda}\times\vec{\mu}_A$ 计算叉积：
\[
\vec{a}\times\vec{\lambda} = a_\mu(\vec{\mu}_A\times\vec{\lambda}) + a_\nu(\vec{\nu}_A\times\vec{\lambda})
                          = -a_\mu\vec{\nu}_A + a_\nu\vec{\mu}_A
\]
即 $a_\mu\vec{\nu}_A - a_\nu\vec{\mu}_A = -\vec{a}\times\vec{\lambda}$，又 $a_\mu\vec{\mu}_A + a_\nu\vec{\nu}_A = \vec{a} - a_\lambda\vec{\lambda}$，代入 $\vec{b}$ 的表达式：
\begin{align*}
    \vec{b} &= (\vec{a}\cdot\vec{\lambda})\vec{\lambda}
             + \bigl[\vec{a} - (\vec{a}\cdot\vec{\lambda})\vec{\lambda}\bigr]\cos\theta
             - (\vec{a}\times\vec{\lambda})\sin\theta \\[4pt]
            &= \vec{a}\cos\theta - \vec{a}\times\vec{\lambda}\sin\theta
             + \vec{\lambda}(\vec{\lambda}\cdot\vec{a})(1-\cos\theta)
\end{align*}
\end{proof}

同样的，对任意一个并矢（Dyadic）也有类似的旋转公式。
\begin{myformula}
    \mathbf{C}  = \mathbf{U} \cos\theta - \mathbf{U} \times\vec{\lambda}\sin\theta + \vec{\lambda}\vec{\lambda}(1-\cos\theta)
\end{myformula}

\section{方向余弦 (Direction Cosines)}

若 \(\mathbf{a}_1, \mathbf{a}_2, \mathbf{a}_3\) 和 \(\mathbf{b}_1, \mathbf{b}_2, \mathbf{b}_3\) 是两组正交单位向量的右手系，那么九个量 \(C_{ij}\)（\(i, j = 1, 2, 3\)），称为\textbf{方向余弦}，定义为：
\begin{equation}
C_{ij} \triangleq \mathbf{a}_i \cdot \mathbf{b}_j \qquad (i, j = 1, 2, 3) \tag{1}
\end{equation}
那么两个行矩阵 \([\mathbf{a}_1 \quad \mathbf{a}_2 \quad \mathbf{a}_3]\) 和 \([\mathbf{b}_1 \quad \mathbf{b}_2 \quad \mathbf{b}_3]\) 之间的关系如下：
\begin{equation}
[\mathbf{b}_1 \quad \mathbf{b}_2 \quad \mathbf{b}_3] = [\mathbf{a}_1 \quad \mathbf{a}_2 \quad \mathbf{a}_3] C \tag{2}
\end{equation}
其中 \(C\) 是一个定义为如下的方阵：
\begin{equation}
C \triangleq \begin{bmatrix} 
C_{11} & C_{12} & C_{13} \\ 
C_{21} & C_{22} & C_{23} \\ 
C_{31} & C_{32} & C_{33} 
\end{bmatrix} \tag{3}
\end{equation}
如果使用上标 \(T\) 表示转置，即 \(C^T\) 定义为：
\begin{equation}
C^T \triangleq \begin{bmatrix} 
C_{11} & C_{21} & C_{31} \\ 
C_{12} & C_{22} & C_{32} \\ 
C_{13} & C_{23} & C_{33} 
\end{bmatrix} \tag{4}
\end{equation}
那么式 (2) 可由等价的如下关系替换：
\begin{equation}
[\mathbf{a}_1 \quad \mathbf{a}_2 \quad \mathbf{a}_3] = [\mathbf{b}_1 \quad \mathbf{b}_2 \quad \mathbf{b}_3] C^T \tag{5}
\end{equation}

矩阵 \(C\) 被称为\textbf{方向余弦矩阵}，可用于描述两个参考系或刚体 \(A\) 和 \(B\) 的相对方位。在那种语境下，使用更丰富的符号 \({}^A C^B\) 替代符号 \(C\) 会更有优势。鉴于式 (2) 和 (5)，必须认为上下标的互换意味着转置；即：
\begin{equation}
{}^B C^A = ({}^A C^B)^T \tag{6}
\end{equation}

方向余弦矩阵 \(C\) 在许多有用的关系中起着重要作用。例如，如果 \(\mathbf{v}\) 是任意向量，\( {}^A v_i\) 和 \( {}^B v_i\)（\(i = 1, 2, 3\)）定义为：
\begin{equation}
{}^A v_i \triangleq \mathbf{v} \cdot \mathbf{a}_i \qquad (i = 1, 2, 3) \tag{7}
\end{equation}
以及：
\begin{equation}
{}^B v_i \triangleq \mathbf{v} \cdot \mathbf{b}_i \qquad (i = 1, 2, 3) \tag{8}
\end{equation}
同时 \({}^A v\) 和 \({}^B v\) 分别表示行矩阵，其元素为 \({}^A v_1, {}^A v_2, {}^A v_3\) 和 \({}^B v_1, {}^B v_2, {}^B v_3\)，那么：
\begin{equation}
{}^B v = {}^A v C \tag{9}
\end{equation}
类似地，如果 \(\mathbf{D}\) 是任意并矢，并且 \({}^A D_{ij}\) 和 \({}^B D_{ij}\)（\(i, j = 1, 2, 3\)）定义为：
\begin{equation}
{}^A D_{ij} \triangleq \mathbf{a}_i \cdot \mathbf{D} \cdot \mathbf{a}_j \qquad (i, j = 1, 2, 3) \tag{10}
\end{equation}
以及：
\begin{equation}
{}^B D_{ij} \triangleq \mathbf{b}_i \cdot \mathbf{D} \cdot \mathbf{b}_j \qquad (i, j = 1, 2, 3) \tag{11}
\end{equation}
而 \({}^A D\) 和 \({}^B D\) 分别表示元素为 \({}^A D_{ij}\) 和 \({}^B D_{ij}\) 的方阵，则：
\begin{equation}
{}^B D = C^T {}^A D C \tag{12}
\end{equation}

对重复下标使用求和约定，通常可以更加简明地表达重要的关系式。例如，\(\delta_{ij}\) 定义为：
\begin{equation}
\delta_{ij} \triangleq 1 - \frac{1}{4}(i-j)^2\left[5-(i-j)^2\right] \qquad (i, j = 1, 2, 3) \tag{13}
\end{equation}
使得当两个下标相同时，\(\delta_{ij}\) 等于 \(1\)，不相同时等于 \(0\)。利用求和约定，可以导出一组控制方向余弦的六个关系式：
\begin{equation}
C_{ik} C_{jk} = \delta_{ij} \qquad (i, j = 1, 2, 3) \tag{14}
\end{equation}
或等价的一组：
\begin{equation}
C_{ki} C_{kj} = \delta_{ij} \qquad (i, j = 1, 2, 3) \tag{15}
\end{equation}
或者，这些关系式可以用矩阵形式表示为：
\begin{equation}
C C^T = U \tag{16}
\end{equation}
和：
\begin{equation}
C^T C = U \tag{17}
\end{equation}
其中 \(U\) 表示单位矩阵，定义为：
\begin{equation}
U \triangleq \begin{bmatrix} 1 & 0 & 0 \\ 0 & 1 & 0 \\ 0 & 0 & 1 \end{bmatrix} \tag{18}
\end{equation}

矩阵 \(C\) 的每个元素均等于其在行列式中的代数余子式；若 \(|C|\) 表示该行列式，则有：
\begin{equation}
|C| = 1 \tag{19}
\end{equation}
因此，\(C\) 是一个正交矩阵，即其逆矩阵等于其转置矩阵。此外：
\begin{equation}
|C - U| = 0 \tag{20}
\end{equation}
因此，1 是每一个方向余弦矩阵的一个特征值。换言之，对于每一个方向余弦矩阵 \(C\)，存在行矩阵 \([\kappa_1 \quad \kappa_2 \quad \kappa_3]\)，称为\textbf{特征向量}，满足方程：
\begin{equation}
[\kappa_1 \quad \kappa_2 \quad \kappa_3] C = [\kappa_1 \quad \kappa_2 \quad \kappa_3] \tag{21}
\end{equation}

现在假定 \(\mathbf{a}_i\) 和 \(\mathbf{b}_i\)（\(i=1, 2, 3\)）分别在参考系或刚体 \(A\) 和 \(B\) 中固定；进一步假定在旋转前 \(\mathbf{a}_i = \mathbf{b}_i\)（\(i=1, 2, 3\)），且 \(\vec{\lambda}\) 和 \(\theta\) 如前文所定义，那么定义 \(\lambda_i\) 为：
\begin{equation}
\lambda_i \triangleq \vec{\lambda} \cdot \mathbf{a}_i = \vec{\lambda} \cdot \mathbf{b}_i \qquad (i = 1, 2, 3) \tag{22}
\end{equation}
那么矩阵 \(C\) 的元素由下式给出：
\begin{align}
C_{11} &= \cos\theta + \lambda_1^2(1-\cos\theta) \tag{23} \\
C_{12} &= -\lambda_3 \sin\theta + \lambda_1\lambda_2(1-\cos\theta) \tag{24} \\
C_{13} &= \lambda_2 \sin\theta + \lambda_3\lambda_1(1-\cos\theta) \tag{25} \\
C_{21} &= \lambda_3 \sin\theta + \lambda_1\lambda_2(1-\cos\theta) \tag{26} \\
C_{22} &= \cos\theta + \lambda_2^2(1-\cos\theta) \tag{27} \\
C_{23} &= -\lambda_1 \sin\theta + \lambda_2\lambda_3(1-\cos\theta) \tag{28} \\
C_{31} &= -\lambda_2 \sin\theta + \lambda_3\lambda_1(1-\cos\theta) \tag{29} \\
C_{32} &= \lambda_1 \sin\theta + \lambda_2\lambda_3(1-\cos\theta) \tag{30} \\
C_{33} &= \cos\theta + \lambda_3^2(1-\cos\theta) \tag{31}
\end{align}
在定义 \(\epsilon_{ijk}\) 之后，方程 (23)--(31) 可以用更简洁的形式表达为：
\begin{equation}
\epsilon_{ijk} \triangleq \frac{1}{2}(i-j)(j-k)(k-i) \qquad (i, j, k = 1, 2, 3) \tag{32}
\end{equation}
（注：当两个或三个下标相同时，\(\epsilon_{ijk}\) 为 0；当下标呈循环顺序时，即 \((1,2,3)\)、\((2,3,1)\) 或 \((3,1,2)\) 时，\(\epsilon_{ijk}\) 为 1；其余情况为 -1）。利用求和约定，式 (23)--(31) 可以替换为：
\begin{equation}
C_{ij} = \delta_{ij} \cos\theta - \epsilon_{ijk}\lambda_k \sin\theta + \lambda_i\lambda_j(1-\cos\theta) \qquad (i, j = 1, 2, 3) \tag{33}
\end{equation}
或者，\(C_{ij}\) 可以用式 (1.1.2) 中定义的并矢 \(\mathbf{C}\) 表示为：
\begin{equation}
C_{ij} = \mathbf{a}_i \cdot (\mathbf{a}_j \cdot \mathbf{C}) \qquad (i, j = 1, 2, 3) \tag{34}
\end{equation}

当 \(\vec{\lambda}\) 平行于 \(\mathbf{a}_i\)，进而平行于 \(\mathbf{b}_i\)（\(i=1, 2, 3\)）时，上述所有结果将大幅简化。如果用 \(C_i(\theta)\) 表示当 \(\vec{\lambda} = \mathbf{a}_i = \mathbf{b}_i\) 时对应的矩阵 \(C\)，则有：
\begin{align}
C_1(\theta) &= \begin{bmatrix} 1 & 0 & 0 \\ 0 & \cos\theta & -\sin\theta \\ 0 & \sin\theta & \cos\theta \end{bmatrix} \tag{35} \\
C_2(\theta) &= \begin{bmatrix} \cos\theta & 0 & \sin\theta \\ 0 & 1 & 0 \\ -\sin\theta & 0 & \cos\theta \end{bmatrix} \tag{36} \\
C_3(\theta) &= \begin{bmatrix} \cos\theta & -\sin\theta & 0 \\ \sin\theta & \cos\theta & 0 \\ 0 & 0 & 1 \end{bmatrix} \tag{37}
\end{align}

前面曾提及，1 是每个方向余弦矩阵的一个特征值。如果矩阵 \(C\) 的元素由式 (23)--(31) 给出，那么行矩阵 \([\lambda_1 \quad \lambda_2 \quad \lambda_3]\) 即为与单位特征值 \(1\) 对应的特征向量，即：
\begin{equation}
[\lambda_1 \quad \lambda_2 \quad \lambda_3] C = [\lambda_1 \quad \lambda_2 \quad \lambda_3] \tag{38}
\end{equation}
等价地，当 \(\mathbf{C}\) 为式 (1.1.2) 中所定义的并矢时：
\begin{equation}
\vec{\lambda} \cdot \mathbf{C} = \vec{\lambda} \tag{39}
\end{equation}

\paragraph{推导部分}
对于任意向量 \(\mathbf{v}\)，以下恒等式成立：
\begin{equation}
\mathbf{v} = (\mathbf{a}_1 \cdot \mathbf{v})\mathbf{a}_1 + (\mathbf{a}_2 \cdot \mathbf{v})\mathbf{a}_2 + (\mathbf{a}_3 \cdot \mathbf{v})\mathbf{a}_3 \tag{40}
\end{equation}
因此，令 \(\mathbf{b}_1\) 扮演 \(\mathbf{v}\) 的角色，可将 \(\mathbf{b}_1\) 表示为：
\begin{equation}
\mathbf{b}_1 = (\mathbf{a}_1 \cdot \mathbf{b}_1)\mathbf{a}_1 + (\mathbf{a}_2 \cdot \mathbf{b}_1)\mathbf{a}_2 + (\mathbf{a}_3 \cdot \mathbf{b}_1)\mathbf{a}_3 \tag{41}
\end{equation}
或者，利用式 (1)，可得：
\begin{equation}
\mathbf{b}_1 = C_{11}\mathbf{a}_1 + C_{21}\mathbf{a}_2 + C_{31}\mathbf{a}_3 \tag{42}
\end{equation}
同理：
\begin{align}
\mathbf{b}_2 &= C_{12}\mathbf{a}_1 + C_{22}\mathbf{a}_2 + C_{32}\mathbf{a}_3 \tag{43} \\
\mathbf{b}_3 &= C_{13}\mathbf{a}_1 + C_{23}\mathbf{a}_2 + C_{33}\mathbf{a}_3 \tag{44}
\end{align}
这三个方程正是基于式 (2) 和矩阵乘法规则形成 \(\mathbf{b}_1, \mathbf{b}_2, \mathbf{b}_3\) 表达式时所得到的结果，类似地通过推理可得到式 (5)。

要验证式 (9) 成立，只需观察：
\begin{equation}
{}^B v_i = \mathbf{v} \cdot (\mathbf{a}_1 C_{1i} + \mathbf{a}_2 C_{2i} + \mathbf{a}_3 C_{3i}) = \mathbf{v} \cdot \mathbf{a}_1 C_{1i} + \mathbf{v} \cdot \mathbf{a}_2 C_{2i} + \mathbf{v} \cdot \mathbf{a}_3 C_{3i} = {}^A v_1 C_{1i} + {}^A v_2 C_{2i} + {}^A v_3 C_{3i} \tag{45}
\end{equation}
并回顾 \({}^A v\) 和 \({}^B v\) 的定义即可。同理，式 (12) 可由下式推导得出：
\begin{align}
{}^B D_{ij} &= (\mathbf{a}_1 C_{1i} + \mathbf{a}_2 C_{2i} + \mathbf{a}_3 C_{3i}) \cdot \mathbf{D} \cdot (\mathbf{a}_1 C_{1j} + \mathbf{a}_2 C_{2j} + \mathbf{a}_3 C_{3j}) \nonumber \\
&= C_{1i} ({}^A D_{11} C_{1j} + {}^A D_{12} C_{2j} + {}^A D_{13} C_{3j}) \nonumber \\
&\quad + C_{2i} ({}^A D_{21} C_{1j} + {}^A D_{22} C_{2j} + {}^A D_{23} C_{3j}) \nonumber \\
&\quad + C_{3i} ({}^A D_{31} C_{1j} + {}^A D_{32} C_{2j} + {}^A D_{33} C_{3j}) \tag{46}
\end{align}
并由 \({}^A D\) 和 \({}^B D\) 的定义得出上述结果。

至于式 (14) 和 (15)，它们是利用求和约定产生的：
\begin{align}
\mathbf{a}_i \cdot \mathbf{a}_j &= C_{ik} C_{jk} \qquad (i, j = 1, 2, 3) \tag{47} \\
\mathbf{b}_i \cdot \mathbf{b}_j &= C_{ki} C_{kj} \qquad (i, j = 1, 2, 3) \tag{48}
\end{align}
分别成立，因为当 \(i = j\) 时 \(\mathbf{a}_i \cdot \mathbf{a}_i\) 等于 1，当 \(i \neq j\) 时等于 0（\(\mathbf{b}_i \cdot \mathbf{b}_j\) 同理）；从而可见式 (16) 和 (17) 是成立的。

此外，将式 (1.1.2) 中给出的并矢 \(\mathbf{C}\) 的表达式代入，可以发现：
\begin{equation}
C_{ij} = \mathbf{a}_i \cdot \mathbf{a}_j \cos\theta - \mathbf{a}_i \cdot \mathbf{a}_j \times \vec{\lambda} \sin\theta + \mathbf{a}_i \cdot \vec{\lambda}\vec{\lambda} \cdot \mathbf{a}_j(1-\cos\theta) \qquad (i, j = 1, 2, 3) \tag{58}
\end{equation}
这与式 (22) 一起可直接推导出式 (23)--(31)；或者结合式 (32)，推导出式 (33)。

可以通过在式 (23)--(31) 中令 \(\lambda_1 = 1\) 且 \(\lambda_2 = \lambda_3 = 0\)，然后利用式 (3) 得到式 (35)。同理，\(\lambda_1 = 0\)、\(\lambda_2 = 1\)、\(\lambda_3 = 0\) 可得到式 (36)，而 \(\lambda_1 = \lambda_2 = 0\)、\(\lambda_3 = 1\) 可得到式 (37)。

最后，式 (39) 由以下观察推导得出：
\begin{align}
\vec{\lambda} \cdot \mathbf{C} &= \vec{\lambda} \cdot \mathbf{U} \cos\theta - \vec{\lambda} \times \vec{\lambda} \sin\theta + \vec{\lambda}^2 \vec{\lambda} (1-\cos\theta) \nonumber \\
&= \vec{\lambda} \cos\theta + 0 + \vec{\lambda} (1-\cos\theta) = \vec{\lambda} \tag{59}
\end{align}
因为 \(\vec{\lambda}\) 是单位向量，故 \(\vec{\lambda}^2 \triangleq \vec{\lambda} \cdot \vec{\lambda} = 1\)；而由式 (22) 和 (34) 即可证明式 (38) 与 (39) 的等价性。

\section{欧拉参数 (Euler Parameters)}

在 1.1 节中引入的单位向量 \(\boldsymbol{\lambda}\) 和角度 \(\theta\)，可用于关联一个向量 \(\boldsymbol{\epsilon}\)（称为**欧拉向量**）以及四个标量 \(\epsilon_1, \dots, \epsilon_4\)（称为**欧拉参数**）到刚体 \(B\) 在参考系 \(A\) 中的一次简单旋转，规定：
\begin{equation}
\boldsymbol{\epsilon} \triangleq \boldsymbol{\lambda} \sin\frac{\theta}{2} \tag{1}
\end{equation}
且：
\begin{equation}
\epsilon_i \triangleq \boldsymbol{\epsilon} \cdot \mathbf{a}_i = \boldsymbol{\epsilon} \cdot \mathbf{b}_i \qquad (i = 1, 2, 3) \tag{2}
\end{equation}
以及：
\begin{equation}
\epsilon_4 \triangleq \cos\frac{\theta}{2} \tag{3}
\end{equation}
其中 \(\mathbf{a}_1, \mathbf{a}_2, \mathbf{a}_3\) 和 \(\mathbf{b}_1, \mathbf{b}_2, \mathbf{b}_3\) 分别为固定在 \(A\) 和 \(B\) 中的单位正交右手向量组，且在旋转前满足 \(\mathbf{a}_i = \mathbf{b}_i\)（\(i=1,2,3\)）。（在涉及两个以上的物体或参考系的讨论中，将使用类似 \({}^A \boldsymbol{\epsilon}^B\) 和 \({}^A\epsilon_i^B\) 的记法。）

欧拉参数并非相互独立的，因为它们的平方和必然等于 1：
\begin{equation}
\epsilon_1^2 + \epsilon_2^2 + \epsilon_3^2 + \epsilon_4^2 = \boldsymbol{\epsilon}^2 + \epsilon_4^2 = 1 \tag{4}
\end{equation}

由 1.2 节引入的方向余弦矩阵 \(C\) 的元素，若用 \(\epsilon_1, \dots, \epsilon_4\) 表示，将呈现出一种特别简单而有规律的形式。若 \(C_{ij}\) 定义为：
\begin{equation}
C_{ij} \triangleq \mathbf{a}_i \cdot \mathbf{b}_j \qquad (i, j = 1, 2, 3) \tag{5}
\end{equation}
则有：
\begin{align}
C_{11} &= \epsilon_1^2 - \epsilon_2^2 - \epsilon_3^2 + \epsilon_4^2 = 1 - 2\epsilon_2^2 - 2\epsilon_3^2 \tag{6} \\
C_{12} &= 2(\epsilon_1 \epsilon_2 - \epsilon_3 \epsilon_4) \tag{7} \\
C_{13} &= 2(\epsilon_3 \epsilon_1 + \epsilon_2 \epsilon_4) \tag{8} \\
C_{21} &= 2(\epsilon_1 \epsilon_2 + \epsilon_3 \epsilon_4) \tag{9} \\
C_{22} &= \epsilon_2^2 - \epsilon_3^2 - \epsilon_1^2 + \epsilon_4^2 = 1 - 2\epsilon_3^2 - 2\epsilon_1^2 \tag{10} \\
C_{23} &= 2(\epsilon_2 \epsilon_3 - \epsilon_1 \epsilon_4) \tag{11} \\
C_{31} &= 2(\epsilon_3 \epsilon_1 - \epsilon_2 \epsilon_4) \tag{12} \\
C_{32} &= 2(\epsilon_2 \epsilon_3 + \epsilon_1 \epsilon_4) \tag{13} \\
C_{33} &= \epsilon_3^2 - \epsilon_1^2 - \epsilon_2^2 + \epsilon_4^2 = 1 - 2\epsilon_1^2 - 2\epsilon_2^2 \tag{14}
\end{align}

欧拉参数也可以用方向余弦来表示，从而使得式 (6)--(14) 能够恒等地成立。这可以通过取以下方式实现：
\begin{align}
\epsilon_1 &= \frac{C_{32} - C_{23}}{4\epsilon_4} \tag{15} \\
\epsilon_2 &= \frac{C_{13} - C_{31}}{4\epsilon_4} \tag{16} \\
\epsilon_3 &= \frac{C_{21} - C_{12}}{4\epsilon_4} \tag{17}
\end{align}
以及：
\begin{equation}
\epsilon_4 = \frac{1}{2} (1 + C_{11} + C_{22} + C_{33})^{1/2} \tag{18}
\end{equation}

如果满足条件：
\begin{equation}
\boldsymbol{\lambda} = \frac{\epsilon_1 \mathbf{a}_1 + \epsilon_2 \mathbf{a}_2 + \epsilon_3 \mathbf{a}_3}{(\epsilon_1^2 + \epsilon_2^2 + \epsilon_3^2)^{1/2}} \tag{19}
\end{equation}
以及：
\begin{equation}
\theta = 2\cos^{-1}\epsilon_4, \qquad 0 \leq \theta \leq \pi \tag{20}
\end{equation}
那么由于式 (1) 和 (3) 成立，我们总可以找到一个简单的旋转，使得与该旋转相关的方向余弦与满足式 (1.2.2) 的任何方向余弦矩阵 \(C\) 的元素对应。换句话说，参考系 \(A\) 和 \(B\) 之间的任意相对方向变化，都可以通过 \(B\) 在 \(A\) 中发生一次简单的旋转来实现。这一命题被称为**欧拉旋转定理 (Euler's theorem on rotation)**。

作为式 (1.1.1) 和 (1.1.3) 的一个替代方案，在参考系 \(A\) 中固定的向量 \(\mathbf{a}\) 与在旋转前与其相等的、固定在刚体 \(B\) 中的向量 \(\mathbf{b}\) 之间的关系，可以用 \(\boldsymbol{\epsilon}\) 和 \(\epsilon_4\) 表示为：
\begin{equation}
\mathbf{b} = \mathbf{a} + 2[\epsilon_4 \boldsymbol{\epsilon} \times \mathbf{a} + \boldsymbol{\epsilon} \times (\boldsymbol{\epsilon} \times \mathbf{a})] \tag{21}
\end{equation}

\paragraph{推导过程}
\(\boldsymbol{\epsilon} \cdot \mathbf{a}_i\) 与 \(\boldsymbol{\epsilon} \cdot \mathbf{b}_i\) 的相等性 [参见式 (2)] 可由式 (1) 和 (1.2.22) 推导得出；式 (4) 则是式 (1)--(3) 的推论，并且利用了 \(\boldsymbol{\lambda}\) 是单位向量这一事实；而式 (6)--(14) 可以通过将式 (1.2.23)--(1.2.31) 中的 \(\theta\) 替换为 \(\theta/2\)，并利用式 (1.2.22) 结合式 (1)--(4) 得到。例如：
\begin{equation}
C_{11} \overset{\text{(1.2.23)}}{=} 2 \cos^2\frac{\theta}{2} - 1 + 2\lambda_1^2 \sin^2\frac{\theta}{2} \tag{22}
\end{equation}
以及：
\begin{equation}
\lambda_1 \sin\frac{\theta}{2} \overset{\text{(1.2.22)}}{=} \boldsymbol{\lambda} \cdot \mathbf{a}_1 \sin\frac{\theta}{2} \overset{\text{(1)}}{=} \boldsymbol{\epsilon} \cdot \mathbf{a}_1 \overset{\text{(2)}}{=} \epsilon_1 \tag{23}
\end{equation}
同时：
\begin{equation}
\cos\frac{\theta}{2} \overset{\text{(3)}}{=} \epsilon_4 \tag{24}
\end{equation}
因此：
\begin{equation}
C_{11} = 2\epsilon_4^2 - 1 + 2\epsilon_1^2 \overset{\text{(4)}}{=} \epsilon_1^2 - \epsilon_2^2 - \epsilon_3^2 + \epsilon_4^2 \tag{25}
\end{equation}
这与式 (6) 一致。

式 (15)--(18) 的有效性可以通过以下方式证明：将式 (15)--(18) 代入式 (6)--(14) 等号右侧，并证明能得到等号左侧的表达式。例如：
\begin{align}
1 - 2\epsilon_2^2 - 2\epsilon_3^2 &= \frac{2(1 + C_{11} + C_{22} + C_{33}) - C_{13}^2 + 2C_{13}C_{31} - C_{31}^2 - C_{21}^2 + 2C_{12}C_{21} - C_{12}^2}{2(1 + C_{11} + C_{22} + C_{33})} \nonumber \\
&= \frac{C_{11} + C_{22} + C_{33} + C_{13}C_{31} + C_{12}C_{21} + C_{11}^2}{1 + C_{11} + C_{22} + C_{33}} \tag{26}
\end{align}
但是，由于矩阵 \(C\) 的每个元素都等于其在 \(|C|\) 中的代数余子式：
\begin{equation}
C_{13}C_{31} = C_{11}C_{33} - C_{22} \tag{27}
\end{equation}
以及：
\begin{equation}
C_{12}C_{21} = C_{11}C_{22} - C_{33} \tag{28}
\end{equation}
因此：
\begin{equation}
1 - 2\epsilon_2^2 - 2\epsilon_3^2 = \frac{C_{11} + C_{11}C_{33} + C_{11}C_{22} + C_{11}^2}{1 + C_{11} + C_{22} + C_{33}} = C_{11} \tag{29}
\end{equation}
正如式 (6) 所要求的。

要证明如果 \(\boldsymbol{\lambda}\) 和 \(\theta\) 分别由式 (19) 和 (20) 给出，则式 (1) 和 (3) 成立，需注意到：
\begin{equation}
\cos\frac{\theta}{2} \overset{\text{(20)}}{=} \epsilon_4 \tag{30}
\end{equation}
这正是式 (3)，并且：
\begin{equation}
\sin\frac{\theta}{2} = (1 - \epsilon_4^2)^{1/2} \overset{\text{(4)}}{=} (\epsilon_1^2 + \epsilon_2^2 + \epsilon_3^2)^{1/2} \tag{31}
\end{equation}
因此：
\begin{equation}
\boldsymbol{\lambda} \sin\frac{\theta}{2} \overset{\text{(19)}}{=} \epsilon_1 \mathbf{a}_1 + \epsilon_2 \mathbf{a}_2 + \epsilon_3 \mathbf{a}_3 \overset{\text{(2)}}{=} \boldsymbol{\epsilon} \tag{32}
\end{equation}
正如式 (1) 所要求的。最后，式 (1.1.1) 等价于：
\begin{align}
\mathbf{b} &= \mathbf{a} + \boldsymbol{\lambda} \times \mathbf{a} \sin\theta + \boldsymbol{\lambda} \times (\boldsymbol{\lambda} \times \mathbf{a})(1 - \cos\theta) \nonumber \\
&\overset{\text{(1)}}{=} \mathbf{a} + 2\boldsymbol{\epsilon} \times \mathbf{a} \cos\frac{\theta}{2} + 2\boldsymbol{\epsilon} \times (\boldsymbol{\epsilon} \times \mathbf{a}) \overset{\text{(3)}}{=} \mathbf{a} + 2[\epsilon_4 \boldsymbol{\epsilon} \times \mathbf{a} + \boldsymbol{\epsilon} \times (\boldsymbol{\epsilon} \times \mathbf{a})] \tag{33}
\end{align}
这与式 (21) 一致。
\section{罗德里格斯参数 (Rodrigues Parameters)}

刚体 \(B\) 在参考系 \(A\) 中的一次**简单旋转**（参见 1.1 节）可以关联一个向量 \(\boldsymbol{\rho}\) 以及三个标量 \(\rho_1, \rho_2, \rho_3\)。其中，\(\boldsymbol{\rho}\) 称为\textbf{罗德里格斯向量 (Rodrigues vector)}，\(\rho_1, \rho_2, \rho_3\) 称为\textbf{罗德里格斯参数 (Rodrigues parameters)}。它们的定义如下：

\begin{mydefinition}{罗德里格斯参数}
\[
\boldsymbol{\rho} \triangleq \boldsymbol{\lambda} \tan \frac{\theta}{2}
\]
\end{mydefinition}
\begin{equation}
\rho_i \triangleq \boldsymbol{\rho} \cdot \mathbf{a}_i = \boldsymbol{\rho} \cdot \mathbf{b}_i \qquad (i = 1, 2, 3) \tag{2}
\end{equation}




其中，\(\boldsymbol{\lambda}\) 和 \(\theta\) 与 1.1 节中的含义相同；\(\mathbf{a}_1, \mathbf{a}_2, \mathbf{a}_3\) 和 \(\mathbf{b}_1, \mathbf{b}_2, \mathbf{b}_3\) 分别是固定在参考系 \(A\) 和刚体 \(B\) 中的右手单位正交向量组，且在旋转前满足 \(\mathbf{a}_i = \mathbf{b}_i \ (i=1,2,3)\)。当讨论涉及两个以上的物体或参考系时，会使用诸如 \({}^A \boldsymbol{\rho}^B\) 和 \({}^A \rho_i^B\) 这样的符号。

\subsection*{与欧拉参数的关系}
罗德里格斯参数与\textbf{欧拉参数}（参见 1.3 节）密切相关，其关系为：
\begin{equation}
\rho_i = \frac{\epsilon_i}{\epsilon_4} \qquad (i = 1, 2, 3) \tag{3}
\end{equation}
罗德里格斯参数相比欧拉参数的主要优点是数量更少。但它的缺点在于，当旋转角度 \(\theta = \pi\) 时，罗德里格斯参数将变为\textbf{无穷大}；而欧拉参数的绝对值则永远不会超过 1。

\subsection*{用罗德里格斯参数表示的方向余弦矩阵 \(C\)}
利用罗德里格斯参数，方向余弦矩阵 \(C\)（参见 1.2 节）可以写成下列紧凑的形式：

\begin{equation}
C = \frac{1}{1 + \rho_1^2 + \rho_2^2 + \rho_3^2} 
\begin{bmatrix} 
1 + \rho_1^2 - \rho_2^2 - \rho_3^2 & 2(\rho_1\rho_2 - \rho_3) & 2(\rho_3\rho_1 + \rho_2) \\
2(\rho_1\rho_2 + \rho_3) & 1 - \rho_1^2 + \rho_2^2 - \rho_3^2 & 2(\rho_2\rho_3 - \rho_1) \\
2(\rho_3\rho_1 - \rho_2) & 2(\rho_2\rho_3 + \rho_1) & 1 - \rho_1^2 - \rho_2^2 + \rho_3^2
\end{bmatrix} \tag{4}
\end{equation}

\subsection*{向量关系与罗德里格斯向量表达式}
罗德里格斯向量可用于建立 1.1 节中定义的旋转前后向量 \(\mathbf{a}\) 和 \(\mathbf{b}\) 的差与和之间的简单关系：
\begin{equation}
\mathbf{a} - \mathbf{b} = (\mathbf{a} + \mathbf{b}) \times \boldsymbol{\rho} \tag{5}
\end{equation}

该关系非常有用。例如，可用它推导出由特定相对方向变化产生的简单旋转所对应的罗德里格斯向量表达式：
\begin{equation}
\boldsymbol{\rho} = \frac{(\boldsymbol{\alpha}_1 - \boldsymbol{\beta}_1) \times (\boldsymbol{\alpha}_2 - \boldsymbol{\beta}_2)}{(\boldsymbol{\alpha}_1 + \boldsymbol{\beta}_1) \cdot (\boldsymbol{\alpha}_2 - \boldsymbol{\beta}_2)} \tag{6}
\end{equation}
其中，\(\boldsymbol{\alpha}_i\) 和 \(\boldsymbol{\beta}_i\)（\(i = 1, 2\)）分别是固定在 \(A\) 和 \(B\) 中的向量，且在相对方向改变前满足 \(\boldsymbol{\alpha}_i = \boldsymbol{\beta}_i\)。

\subsection*{由方向余弦求罗德里格斯参数}
反过来，给定方向余弦矩阵，其对应的罗德里格斯参数可通过下式求出：
\begin{align}
\rho_1 &= \frac{C_{32} - C_{23}}{1 + C_{11} + C_{22} + C_{33}} \tag{7} \\
\rho_2 &= \frac{C_{13} - C_{31}}{1 + C_{11} + C_{22} + C_{33}} \tag{8} \\
\rho_3 &= \frac{C_{21} - C_{12}}{1 + C_{11} + C_{22} + C_{33}} \tag{9}
\end{align}

\subsection*{推导部分 (Derivations)}
\begin{itemize}
    \item \textbf{式 (2) 的证明}：由式 (1) 和式 (1.2.22) 成立。
    \item \textbf{式 (3) 的证明}：因为 \(\frac{\boldsymbol{\epsilon}}{\epsilon_4} = \boldsymbol{\lambda} \tan \frac{\theta}{2} = \boldsymbol{\rho}\)，所以 \(\rho_i = \frac{\boldsymbol{\epsilon} \cdot \mathbf{a}_i}{\epsilon_4} = \frac{\epsilon_i}{\epsilon_4}\)。
    \item \textbf{方向余弦矩阵的推导}：由欧拉参数平方和为 1 的性质 \([1 = (\rho_1^2 + \rho_2^2 + \rho_3^2 + 1)\epsilon_4^2]\) 得到 \(\epsilon_4^2\) 的表达式，并结合欧拉参数表示的方向余弦矩阵（式 1.3.6-1.3.14），可推得式 (4)。
    \item \textbf{向量关系式 (5) 的证明}：对式 (1.3.21) 两边叉乘 \(\boldsymbol{\epsilon}\) 并利用向量三重积恒等式化简，最终得到 \(\mathbf{a} - \mathbf{b} = (\mathbf{a} + \mathbf{b}) \times \frac{\boldsymbol{\epsilon}}{\epsilon_4} = (\mathbf{a} + \mathbf{b}) \times \boldsymbol{\rho}\)。
    \item \textbf{式 (6)-(9) 的证明}：利用上述向量关系进行叉积运算，并结合 \(\rho_1 = \frac{\epsilon_1}{\epsilon_4} = \frac{C_{32} - C_{23}}{4\epsilon_4^2} = \frac{C_{32} - C_{23}}{1 + C_{11} + C_{22} + C_{33}}\) 即可得到。
\end{itemize}



\section{间接确定方向 (Indirect Determination of Orientation)}

如果 \(\mathbf{a}_1, \mathbf{a}_2, \mathbf{a}_3\) 和 \(\mathbf{b}_1, \mathbf{b}_2, \mathbf{b}_3\) 分别是固定在参考系 \(A\) 和 \(B\) 中的右手单位正交向量组，且每个单位向量在两个参考系中的方向均为已知，那么两个参考系的相对方向可以通过方向余弦 \(C_{ij} \ (i,j = 1,2,3)\) 来描述，依据式 (1.2.1)，只需计算 \(\mathbf{a}_i \cdot \mathbf{b}_j \ (i,j = 1,2,3)\) 即可。然而，即使这些点积无法如此直接计算，我们也有可能求出 \(C_{ij} \ (i,j = 1,2,3)\)。

例如，当两个不平行向量（记为 \(\mathbf{p}\) 和 \(\mathbf{q}\)）在 \(A\) 和 \(B\) 中的方向均已知时，就可以直接计算点积 \(\mathbf{a}_i \cdot \mathbf{p}, \mathbf{a}_i \cdot \mathbf{q}, \mathbf{b}_i \cdot \mathbf{p}\) 和 \(\mathbf{b}_i \cdot \mathbf{q} \ (i=1,2,3)\)。在这种情况下，可以按如下方法求出 \(C_{ij}\)：构造一个向量 \(\mathbf{r}\) 和一个并矢 \(\boldsymbol{\sigma}\)：
\begin{equation}
\mathbf{r} \triangleq \mathbf{p} \times \mathbf{q} \tag{1}
\end{equation}
\begin{equation}
\boldsymbol{\sigma} \triangleq \frac{\mathbf{p} \times \mathbf{q} \, \mathbf{r} + \mathbf{q} \times \mathbf{r} \, \mathbf{p} + \mathbf{r} \times \mathbf{p} \, \mathbf{q}}{\mathbf{r}^2} \tag{2}
\end{equation}
接下来，将式 (2) 中每个并矢的前项用 \(\mathbf{a}_i\) 表示，后项用 \(\mathbf{b}_i \ (i=1,2,3)\) 表示。最后，执行以下关系式中指示的乘法运算：
\begin{equation}
C_{ij} = \mathbf{a}_i \cdot \boldsymbol{\sigma} \cdot \mathbf{b}_j \qquad (i,j = 1,2,3) \tag{3}
\end{equation}

\subsection*{推导}
可以证明式 (2) 中定义的并矢 \(\boldsymbol{\sigma}\) 是一个单位并矢，即对于任意向量 \(\mathbf{v}\)：
\begin{equation}
\mathbf{v} = \mathbf{v} \cdot \boldsymbol{\sigma} \tag{4}
\end{equation}
那么式 (3) 的有效性便可看作是式 (1.2.1) 中 \(C_{ij}\) 定义的直接推论。

如果 \(\mathbf{p}\) 和 \(\mathbf{q}\) 不平行，且 \(\mathbf{r}\) 如式 (1) 所示定义，那么任意向量 \(\mathbf{v}\) 都可以表示为：
\begin{equation}
\mathbf{v} = \alpha \mathbf{p} + \beta \mathbf{q} + \gamma \mathbf{r} \tag{5}
\end{equation}
其中 \(\alpha, \beta, \gamma\) 为标量。由式 (5)，可以通过分别与 \(\mathbf{q} \times \mathbf{r}\)、\(\mathbf{r} \times \mathbf{p}\) 和 \(\mathbf{p} \times \mathbf{q}\) 进行标量积运算，反解出 \(\alpha, \beta, \gamma\)。将结果代回式 (5)，即可证明 \(\boldsymbol{\sigma}\) 为单位并矢。

\subsection*{例题}
从两个航天器 \(A\) 和 \(B\) 上同时观测两颗恒星 \(P\) 和 \(Q\)，以生成用于确定 \(A\) 和 \(B\) 相对方向的数据。观测包括测量图 1.5.1 中所示的角度 \(\phi\) 和 \(\psi\)。其中 \(O\) 代表 \(A\) 或 \(B\) 中固定的一个点，\(R\) 为 \(P\) 或 \(Q\)，\(\mathbf{c}_1, \mathbf{c}_2, \mathbf{c}_3\) 为固定在 \(A\) 或 \(B\) 中的一组右手单位正交向量。

对于表 1.5.1 中给出的角度数值，需要确定方向余弦矩阵 \(C\)。

\begin{table}[h]
\centering
\caption{表 1.5.1 角度 \(\phi\) 和 \(\psi\)（单位：度）}
\begin{tabular}{l c c c c}
\toprule
 & \multicolumn{2}{c}{\(P\)} & \multicolumn{2}{c}{\(Q\)} \\
\cmidrule(lr){2-3} \cmidrule(lr){4-5}
 & \(\phi\) & \(\psi\) & \(\phi\) & \(\psi\) \\
\midrule
\(\mathbf{c}_i = \mathbf{a}_i\) & 90 & 45 & 30 & 0 \\
\(\mathbf{c}_i = \mathbf{b}_i\) & 135 & 0 & 90 & 60 \\
\bottomrule
\end{tabular}
\end{table}

\begin{table}[h]
\centering
\caption{表 1.5.2 出现在 \(\boldsymbol{\sigma}\) 中的向量}
\resizebox{\textwidth}{!}{%
\begin{tabular}{l c c c c c c c}
\toprule
行 & 向量 & \(\mathbf{a}_1\) & \(\mathbf{a}_2\) & \(\mathbf{a}_3\) & \(\mathbf{b}_1\) & \(\mathbf{b}_2\) & \(\mathbf{b}_3\) \\
\midrule
1 & \(\mathbf{p}\) & \(0\) & \(\frac{\sqrt{2}}{2}\) & \(\frac{\sqrt{2}}{2}\) & \(-\frac{\sqrt{2}}{2}\) & \(\frac{\sqrt{2}}{2}\) & \(0\) \\
2 & \(\mathbf{q}\) & \(\frac{\sqrt{3}}{2}\) & \(\frac{1}{2}\) & \(0\) & \(0\) & \(\frac{1}{2}\) & \(\frac{\sqrt{3}}{2}\) \\
3 & \(\mathbf{r} = \mathbf{p} \times \mathbf{q}\) & \(-\frac{\sqrt{2}}{4}\) & \(\frac{\sqrt{6}}{4}\) & \(-\frac{\sqrt{6}}{4}\) & \(\frac{\sqrt{6}}{4}\) & \(\frac{\sqrt{6}}{4}\) & \(-\frac{\sqrt{2}}{4}\) \\
4 & \(\mathbf{q} \times \mathbf{r}\) & \(-\frac{\sqrt{6}}{8}\) & \(\frac{3\sqrt{2}}{8}\) & \(\frac{\sqrt{2}}{2}\) & \(-\frac{\sqrt{2}}{2}\) & \(\frac{\sqrt{2}}{2}\) & \(0\) \\
5 & \(\mathbf{r} \times \mathbf{p}\) & \(\frac{\sqrt{3}}{2}\) & \(\frac{1}{4}\) & \(-\frac{1}{4}\) & \(0\) & \(\frac{1}{2}\) & \(\frac{\sqrt{3}}{2}\) \\
\bottomrule
\end{tabular}
}

\end{table}

由此可得：
\begin{align}
\boldsymbol{\sigma} &= \left[ \left( -\frac{\sqrt{2}}{4} \mathbf{a}_1 + \frac{\sqrt{6}}{4} \mathbf{a}_2 - \frac{\sqrt{6}}{4} \mathbf{a}_3 \right) \left( \frac{\sqrt{6}}{4} \mathbf{b}_1 + \frac{\sqrt{6}}{4} \mathbf{b}_2 - \frac{\sqrt{2}}{4} \mathbf{b}_3 \right) \right. \nonumber \\
&\quad + \left( -\frac{\sqrt{6}}{8} \mathbf{a}_1 + \frac{3\sqrt{2}}{8} \mathbf{a}_2 + \frac{\sqrt{2}}{2} \mathbf{a}_3 \right) \left( -\frac{\sqrt{2}}{2} \mathbf{b}_1 + \frac{\sqrt{2}}{2} \mathbf{b}_2 + 0\mathbf{b}_3 \right) \nonumber \\
&\quad \left. + \left( \frac{\sqrt{3}}{2} \mathbf{a}_1 + \frac{1}{4} \mathbf{a}_2 - \frac{1}{4} \mathbf{a}_3 \right) \left( 0\mathbf{b}_1 + \frac{1}{2} \mathbf{b}_2 + \frac{\sqrt{3}}{2} \mathbf{b}_3 \right) \right] \Bigg/ \frac{7}{8} \tag{13}
\end{align}
且：
\begin{align}
C_{11} &= \mathbf{a}_1 \cdot \boldsymbol{\sigma} \cdot \mathbf{b}_1 = \left( -\frac{\sqrt{2}}{4} \frac{\sqrt{6}}{4} + \frac{\sqrt{6}}{8} \frac{\sqrt{2}}{2} \right) \Bigg/ \frac{7}{8} = 0 \tag{14} \\
C_{12} &= \mathbf{a}_1 \cdot \boldsymbol{\sigma} \cdot \mathbf{b}_2 = \left( -\frac{\sqrt{2}}{4} \frac{\sqrt{6}}{4} - \frac{\sqrt{6}}{8} \frac{\sqrt{2}}{2} + \frac{\sqrt{3}}{2} \frac{1}{2} \right) \Bigg/ \frac{7}{8} = 0 \tag{15} \\
C_{13} &= \mathbf{a}_1 \cdot \boldsymbol{\sigma} \cdot \mathbf{b}_3 = \left( \frac{\sqrt{2}}{4} \frac{\sqrt{2}}{4} + \frac{\sqrt{3}}{2} \frac{\sqrt{3}}{2} \right) \Bigg/ \frac{7}{8} = 1 \tag{16}
\end{align}
以此类推，得出方向余弦矩阵：
\begin{equation}
C = \begin{bmatrix} 0 & 0 & 1 \\ 0 & 1 & 0 \\ -1 & 0 & 0 \end{bmatrix} \tag{17}
\end{equation}


\section{连续旋转 (Successive Rotations)}

\begin{mydefinition}{连续旋转的定义}
    当刚体 \(B\) 在参考系 \(A\) 中经历两次连续的简单旋转时（参见 1.1 节），每次旋转以及一个等效的单次旋转都可以用方向余弦、欧拉参数和罗德里格斯参数来描述。为了探讨这些关系，引入一个假想的刚体 \(\bar{B}\)。
    \begin{enumerate}
        \item \(\bar{B}\) 在第一次旋转时运动方式与 \(B\) 完全相同；
        \item 在 \(B\) 进行第二次旋转的过程中，\(\bar{B}\) 保持固定于参考系 \(A\) 中。
    \end{enumerate}
    从而，第一次旋转可视为 \(B\) 相对于 \(A\) 的旋转，第二次旋转视为 \(B\) 相对于 \(\bar{B}\) 的旋转。
\end{mydefinition}

设 \(\mathbf{a}_i, \mathbf{b}_i, \bar{\mathbf{b}}_i\ (i=1,2,3)\) 分别为固定在参考系 \(A\)、刚体 \(B\) 和假想刚体 \(\bar{B}\) 中的右手单位正交向量组，且在 \(B\) 相对 \(A\) 第一次旋转前满足 \(\mathbf{a}_i = \mathbf{b}_i = \bar{\mathbf{b}}_i\ (i=1,2,3)\)。
记 \({}^A C^{\bar{B}}\)、\({}^{\bar{B}} C^B\) 和 \({}^A C^B\) 分别表征第一次、第二次以及**等效单次旋转**的**方向余弦矩阵**，使得：
\begin{align}
[\bar{\mathbf{b}}_1 \quad \bar{\mathbf{b}}_2 \quad \bar{\mathbf{b}}_3] &= [\mathbf{a}_1 \quad \mathbf{a}_2 \quad \mathbf{a}_3] {}^A C^{\bar{B}} \tag{1} \\
[\mathbf{b}_1 \quad \mathbf{b}_2 \quad \mathbf{b}_3] &= [\bar{\mathbf{b}}_1 \quad \bar{\mathbf{b}}_2 \quad \bar{\mathbf{b}}_3] {}^{\bar{B}} C^B \tag{2} \\
[\mathbf{b}_1 \quad \mathbf{b}_2 \quad \mathbf{b}_3] &= [\mathbf{a}_1 \quad \mathbf{a}_2 \quad \mathbf{a}_3] {}^A C^B \tag{3}
\end{align}

\begin{myTheorem}{方向余弦矩阵的合成定理}
    连续两次旋转的等效方向余弦矩阵，等于各次旋转方向余弦矩阵的**乘积**：
    \begin{equation}
    {}^A C^B = {}^A C^{\bar{B}} \ {}^{\bar{B}} C^B \tag{4}
    \end{equation}
    \textbf{注意}：由于矩阵乘法不满足交换律，旋转的**顺序不可颠倒**。
\end{myTheorem}

类似地，对于**罗德里格斯向量**（见 1.4 节），设 \({}^A \boldsymbol{\rho}^{\bar{B}}\)、\({}^{\bar{B}} \boldsymbol{\rho}^{B}\) 和 \({}^A \boldsymbol{\rho}^{B}\) 对应各次旋转，则有：
\begin{myformula}[title={罗德里格斯向量合成公式}]
    {}^A \boldsymbol{\rho}^B = \frac{{}^A \boldsymbol{\rho}^{\bar{B}} + {}^{\bar{B}} \boldsymbol{\rho}^B + {}^{\bar{B}} \boldsymbol{\rho}^B \times {}^A \boldsymbol{\rho}^{\bar{B}}}{1 - {}^A \boldsymbol{\rho}^{\bar{B}} \cdot {}^{\bar{B}} \boldsymbol{\rho}^B}
\end{myformula}

若用**欧拉参数**描述，设第一次、第二次及等效单次旋转的欧拉参数分别为 \({}^A \epsilon_i^{\bar{B}}, {}^{\bar{B}} \epsilon_i^{B}, {}^A \epsilon_i^{B}\ (i=1,\dots,4)\)，并满足：
\begin{myformula}[title={欧拉参数合成矩阵}]
    \begin{bmatrix} {}^A \epsilon_1^B \\ {}^A \epsilon_2^B \\ {}^A \epsilon_3^B \\ {}^A \epsilon_4^B \end{bmatrix} = \begin{bmatrix}
    {}^A \epsilon_4^{\bar{B}} & -{}^A \epsilon_3^{\bar{B}} & {}^A \epsilon_2^{\bar{B}} & {}^A \epsilon_1^{\bar{B}} \\
    {}^A \epsilon_3^{\bar{B}} & {}^A \epsilon_4^{\bar{B}} & -{}^A \epsilon_1^{\bar{B}} & {}^A \epsilon_2^{\bar{B}} \\
    -{}^A \epsilon_2^{\bar{B}} & {}^A \epsilon_1^{\bar{B}} & {}^A \epsilon_4^{\bar{B}} & {}^A \epsilon_3^{\bar{B}} \\
    -{}^A \epsilon_1^{\bar{B}} & -{}^A \epsilon_2^{\bar{B}} & -{}^A \epsilon_3^{\bar{B}} & {}^A \epsilon_4^{\bar{B}}
    \end{bmatrix} \begin{bmatrix} {}^{\bar{B}} \epsilon_1^B \\ {}^{\bar{B}} \epsilon_2^B \\ {}^{\bar{B}} \epsilon_3^B \\ {}^{\bar{B}} \epsilon_4^B \end{bmatrix} \\
    \text{等效形式为：}\\
    {}^A \boldsymbol{\epsilon}^B = {}^A \boldsymbol{\epsilon}^{\bar{B}} {}^{\bar{B}} \epsilon_4^B + {}^{\bar{B}} \boldsymbol{\epsilon}^B {}^A \epsilon_4^{\bar{B}} + {}^{\bar{B}} \boldsymbol{\epsilon}^B \times {}^A \boldsymbol{\epsilon}^{\bar{B}} \\
    {}^A \epsilon_4^B = {}^A \epsilon_4^{\bar{B}} {}^{\bar{B}} \epsilon_4^B - {}^A \boldsymbol{\epsilon}^{\bar{B}} \cdot {}^{\bar{B}} \boldsymbol{\epsilon}^B
\end{myformula}

\begin{proof}[合成公式的推导证明]
    将式 (1) 代入式 (2) 得 \([\mathbf{b}_1 \ \mathbf{b}_2 \ \mathbf{b}_3] = [\mathbf{a}_1 \ \mathbf{a}_2 \ \mathbf{a}_3] {}^A C^{\bar{B}} {}^{\bar{B}} C^B\)，与式 (3) 对比即可证明式 (4)。对于罗德里格斯向量的合成，设固定于 \(A\)、\(\bar{B}\) 和 \(B\) 中的向量 \(\mathbf{a}, \bar{\mathbf{b}}, \mathbf{b}\) 满足 \(\mathbf{a}=\bar{\mathbf{b}}=\mathbf{b}\)，通过式 (1.4.5) \(\mathbf{a}-\bar{\mathbf{b}}=(\mathbf{a}+\bar{\mathbf{b}})\times {}^A \boldsymbol{\rho}^{\bar{B}}\) 与 \(\bar{\mathbf{b}}-\mathbf{b}=(\bar{\mathbf{b}}+\mathbf{b})\times {}^{\bar{B}} \boldsymbol{\rho}^{B}\) 联立消去 \(\bar{\mathbf{b}}\)，即可导出式 (5)。式 (12) 至 (14) 则是通过将欧拉参数与罗德里格斯向量的关系 \(\boldsymbol{\rho} = \boldsymbol{\epsilon}/\epsilon_4\) 代入式 (5) 结合恒等式化简得出。
\end{proof}

\begin{example}[陀螺仪支架的连续旋转]
    如图 1.6.1 所示，刚体 \(B\) 首先绕 \(X\) 轴（垂直于 \(\mathbf{a}_1\) 且与 \(\mathbf{a}_2,\mathbf{a}_3\) 成固定角）旋转 \(90^\circ\)，然后绕 \(Y\) 轴（垂直于 \(\mathbf{a}_2\) 且与 \(\mathbf{a}_1,\mathbf{a}_3\) 成固定角）旋转 \(180^\circ\)。求先后发生这两次旋转（即先 \(X\) 后 \(Y\)）时的等效方向余弦矩阵 \(C_x\)，以及先 \(Y\) 后 \(X\) 时的等效矩阵 \(C_y\)。
    \begin{solution}
    对于先 \(X\) 后 \(Y\)，先求出第一次旋转矩阵 \({}^A C^{\bar{B}}\)，再求出第二次旋转矩阵 \({}^{\bar{B}} C^B\)，利用式 (4) 相乘得：
    \begin{equation}
        C_x = {}^A C^{\bar{B}} {}^{\bar{B}} C^B = \begin{bmatrix} 0 & \frac{1}{2} & -\frac{\sqrt{3}}{2} \\ -1 & 0 & 0 \\ 0 & \frac{\sqrt{3}}{2} & \frac{1}{2} \end{bmatrix} \tag{41}
    \end{equation}
    另一种解法是使用欧拉参数。根据题意 \(\theta_1=\pi, \lambda=\frac{1}{2}\mathbf{a}_2+\frac{\sqrt{3}}{2}\mathbf{a}_3\)，求得第一次旋转的欧拉参数，带入矩阵合成公式 (12)，再转换为方向余弦矩阵，同样可求得 \(C_x\)。
    同理可得先 \(Y\) 后 \(X\) 的矩阵 \(C_y\) 为：
    \begin{equation}
        C_y = \begin{bmatrix} 0 & 1 & 0 \\ -\frac{1}{2} & 0 & \frac{\sqrt{3}}{2} \\ \frac{\sqrt{3}}{2} & 0 & \frac{1}{2} \end{bmatrix} \tag{49}
    \end{equation}
    可见两种顺序产生的结果完全不同，再次印证了旋转的**非交换性**。
    \end{solution}
\end{example}

\vspace{1em}
\section{方向角 (Orientation Angles)}

\begin{mydefinition}{描述方向的四种角度系统}
    刚体 \(B\) 在参考系 \(A\) 中的方向有时需要用三个角度来描述（称为**方向角**）。
    根据旋转轴是固定在参考系 \(A\) 中，还是固定在刚体 \(B\) 中，可以分为：
    \begin{itemize}
        \item \textbf{空间-三角度 (Space-three angles)}：连续绕 \(A\) 系中固定轴 \(\mathbf{a}_1, \mathbf{a}_2, \mathbf{a}_3\) 旋转；
        \item \textbf{体-三角度 (Body-three angles)}：连续绕 \(B\) 系中固定轴 \(\mathbf{b}_1, \mathbf{b}_2, \mathbf{b}_3\) 旋转；
        \item \textbf{空间-二角度 (Space-two angles)}：连续绕 \(\mathbf{a}_1, \mathbf{a}_2, \mathbf{a}_1\) 旋转；
        \item \textbf{体-二角度 (Body-two angles)}：连续绕 \(\mathbf{b}_1, \mathbf{b}_2, \mathbf{b}_1\) 旋转。
    \end{itemize}
\end{mydefinition}

设旋转量分别为 \(\theta_1, \theta_2, \theta_3\)，记 \(s_i = \sin\theta_i, c_i=\cos\theta_i\ (i=1,2,3)\)。

\begin{myformula}[title={空间-三角度方向余弦矩阵}]
    C = \begin{bmatrix} c_2 c_3 & s_1 s_2 c_3 - s_3 c_1 & c_1 s_2 c_3 + s_3 s_1 \\ c_2 s_3 & s_1 s_2 s_3 + c_3 c_1 & c_1 s_2 s_3 - c_3 s_1 \\ -s_2 & s_1 c_2 & c_1 c_2 \end{bmatrix}
\end{myformula}
注意，当 \(|C_{31}| = 1\) 时会出现奇异性（通常称为万向节死锁），需要特殊处理。
\begin{myformula}[title={体-三角度方向余弦矩阵}]
    C = \begin{bmatrix} c_2 c_3 & -c_2 s_3 & s_2 \\ s_1 s_2 c_3 + s_3 c_1 & -s_1 s_2 s_3 + c_3 c_1 & -s_1 c_2 \\ -c_1 s_2 c_3 + s_3 s_1 & c_1 s_2 s_3 + c_3 s_1 & c_1 c_2 \end{bmatrix}
\end{myformula}

\begin{myformula}[title={空间-二角度方向余弦矩阵}]
    C = \begin{bmatrix} c_2 & s_1 s_2 & c_1 s_2 \\ s_2 s_3 & -s_1 c_2 s_3 + c_3 c_1 & -c_1 c_2 s_3 - c_3 s_1 \\ -s_2 c_3 & s_1 c_2 c_3 + s_3 c_1 & c_1 c_2 c_3 - s_3 s_1 \end{bmatrix}
\end{myformula}

\begin{myformula}[title={体-二角度方向余弦矩阵}]
    C = \begin{bmatrix} c_2 & s_2 s_3 & s_2 c_3 \\ s_1 s_2 & -s_1 c_2 s_3 + c_3 c_1 & -s_1 c_2 c_3 - s_3 c_1 \\ -c_1 s_2 & c_1 c_2 s_3 + c_3 s_1 & c_1 c_2 c_3 - s_3 s_1 \end{bmatrix}
\end{myformula}

\begin{proof}[体三矩阵的推导]
    以体-三角度为例，利用式 (1.6.15) 三次旋转合成定理 \({}^A C^B = {}^A C^{\bar{B}} \ {}^{\bar{B}} C^{\bar{\bar{B}}} \ {}^{\bar{\bar{B}}} C^B\)，令：
    \begin{equation*}
        {}^A C^{\bar{B}} = \begin{bmatrix} 1 & 0 & 0 \\ 0 & c_1 & -s_1 \\ 0 & s_1 & c_1 \end{bmatrix}, \quad
        {}^{\bar{B}} C^{\bar{\bar{B}}} = \begin{bmatrix} c_2 & 0 & s_2 \\ 0 & 1 & 0 \\ -s_2 & 0 & c_2 \end{bmatrix}, \quad
        {}^{\bar{\bar{B}}} C^B = \begin{bmatrix} c_3 & -s_3 & 0 \\ s_3 & c_3 & 0 \\ 0 & 0 & 1 \end{bmatrix}
    \end{equation*}
    将上述三个矩阵依次相乘，即可推导出式 (11)。空间-二、体-二角度的矩阵证明过程类似。
\end{proof}

\begin{example}[陀螺仪系统的角度转换]
    如图 1.7.2 所示，某陀螺仪系统，其外框架轴（\(\mathbf{a}_1\)）、内框架轴以及转子轴（\(\mathbf{b}_1\)）使用了**体-二角度** \(\theta_1, \theta_2, \theta_3\) 描述（即 \(\mathbf{b}_1, \mathbf{b}_2, \mathbf{b}_1\) 轴旋转）。
    在某个时刻，\(\theta_1 = 30^\circ, \theta_2 = 45^\circ, \theta_3 = 60^\circ\)，且对应角度变化率为 \(\dot{\theta}_1=1.00, \dot{\theta}_2=2.00, \dot{\theta}_3=3.00\) rad/sec。试求此时对应的**空间-三角度** \(\phi_1, \phi_2, \phi_3\) 及其角速度变化率 \(\dot{\phi}_1, \dot{\phi}_2, \dot{\phi}_3\)。
    \begin{solution}
    首先由体-二矩阵 (31) 计算出方向余弦矩阵的各项值：\(C_{31} = -0.612, C_{32} = 0.780, C_{33} = -0.127, C_{21} = 0.354, C_{11} = 0.707\)。
    将 \(C_{31}\) 代入式 (1.7.2)~(1.7.6)（空间-三角度的反解公式）中，可求得：
    \[
    \phi_2 = 37.7^\circ, \quad \phi_1 = 99.6^\circ, \quad \phi_3 = 26.6^\circ
    \]
    接着利用空间三角度的角速度反解公式 (1.17.2) 和 (1.17.4)，计算出 \(\omega_1, \omega_2, \omega_3\) 后，代入对应公式可得：
    \[
    \dot{\phi}_1 = 5.12 \ \text{rad/sec}, \quad \dot{\phi}_2 = 1.08 \ \text{rad/sec}, \quad \dot{\phi}_3 = 2.31 \ \text{rad/sec}
    \]
    \end{solution}
\end{example}
\section{小旋转 (Small Rotations)}

\begin{mydefinition}{小旋转的近似}
当简单旋转（1.1节）的角度 \(\theta\) 足够小，以至于在分析中可以忽略 \(\theta\) 的二阶及更高阶项时，之前讨论的许多关系可以用更简单的形式替代。
\end{mydefinition}

在忽略 \(\theta\) 的高阶项后，有如下近似：
\begin{align}
\mathbf{b} &= \mathbf{a} - \mathbf{a} \times \boldsymbol{\lambda} \theta \tag{1} \\
\mathbf{C} &= \mathbf{U} - \mathbf{U} \times \boldsymbol{\lambda} \theta \tag{2} \\
\boldsymbol{\epsilon} &= \frac{1}{2} \boldsymbol{\lambda} \theta \tag{3} \\
\epsilon_4 &= 1 \tag{4} \\
\boldsymbol{\rho} &= \frac{1}{2} \boldsymbol{\lambda} \theta \tag{5}
\end{align}
这表明，在该近似精度下，罗德里格斯向量与欧拉向量是不可区分的。

\begin{mydefinition}{反对称矩阵记号}
在处理小旋转时，经常涉及反对称矩阵。为了方便，建立如下记号：在某个字母（如 \(q\)）上方加波浪线 \(\widetilde{q}\)，表示一个反对称矩阵，其非对角元素由 \(\pm q_i\ (i=1,2,3)\) 组成：
\begin{equation}
\widetilde{q} = \begin{bmatrix} 0 & -q_3 & q_2 \\ q_3 & 0 & -q_1 \\ -q_2 & q_1 & 0 \end{bmatrix} \tag{6}
\end{equation}
\end{mydefinition}

利用此约定，对小旋转有如下结果：
\begin{equation}
C = U + \widetilde{\boldsymbol{\lambda}} \theta \tag{7}
\end{equation}
由欧拉参数表示，则为：
\begin{equation}
C = U + 2\widetilde{\boldsymbol{\epsilon}} \tag{8}
\end{equation}
若考虑两次连续的小旋转，单次等效小旋转的方向余弦矩阵为：
\begin{equation}
{}^A C^B = U + ({}^A C^{\bar{B}} + {}^{\bar{B}} C^B - 2U) \tag{9}
\end{equation}
对于三次连续旋转，有：
\begin{equation}
{}^A C^B = U + ({}^A C^{\bar{B}} + {}^{\bar{B}} C^{\bar{\bar{B}}} + {}^{\bar{\bar{B}}} C^B - 3U) \tag{10}
\end{equation}
若用罗德里格斯向量表示，则满足简单的加法：
\begin{equation}
{}^A \boldsymbol{\rho}^B = {}^A \boldsymbol{\rho}^{\bar{B}} + {}^{\bar{B}} \boldsymbol{\rho}^B \tag{11}
\end{equation}
对于方向角 \(\theta_1, \theta_2, \theta_3\) 足够小的情况，空间三角度和体三角度的方向余弦矩阵均退化为：
\begin{equation}
C = U + \widetilde{\boldsymbol{\theta}} \tag{12}
\end{equation}
其中 \(\widetilde{\boldsymbol{\theta}}\) 为 \(\theta_i\) 对应的反对称矩阵。

\begin{proof}[近似公式推导]
式 (1) 和 (2) 分别由 (1.1.1) 和 (1.1.2) 将 \(\sin\theta\) 替换为 \(\theta\)，\(\cos\theta\) 替换为 \(1\) 得到。式 (3)-(5) 同理。
式 (9) 的推导是将 \({}^A C^{\bar{B}} = U + \widetilde{\boldsymbol{\lambda}}\theta\) 和 \({}^{\bar{B}} C^B = U + \widetilde{\boldsymbol{\mu}}\phi\) 代入 (1.6.4)，并舍去二阶乘积项得到。式 (11) 可由 (1.6.5) 结合 (5) 推导得出。
\end{proof}

\begin{example}[薄板的小转动位移]
如图 1.8.1 所示，矩形薄板 \(B\) 上一点 \(P\) 依次绕 \(X\) 轴转动 \(0.01\) rad，绕 \(Y\) 轴转动 \(0.02\) rad。假定这两个角度足够小，求 \(P\) 点移动的距离 \(d\)。
记转动前 \(P\) 的位矢为 \(\mathbf{a}=L(3\mathbf{a}_1+4\mathbf{a}_2)\)。等效转轴 \(\boldsymbol{\lambda}\) 通过合成公式 (11) 得到 \(\lambda\theta = 0.01\boldsymbol{\lambda}_x + 0.02\boldsymbol{\lambda}_y\)，代入 (1) 求得位移，其大小为 \(d \approx 0.098L\)。
\end{example}

\section{螺旋运动 (Screw Motion)}

\begin{mydefinition}{刚体位移与螺旋运动}
若刚体 \(B\) 上所有点的位移矢量相等，则称之为**平移**。刚体 \(B\) 在参考系 \(A\) 中的任何位移都可以通过一次平移（将任意基点 \(P\) 移至终点）和一次简单旋转（此时 \(P\) 保持不动）来产生。当基点的位移矢量平行于旋转的罗德里格斯向量时，这种位移称为**螺旋运动**。
任何刚体位移都可以由螺旋运动产生。满足此条件的基点全部位于一条平行于罗德里格斯向量的直线上，这条线称为**螺旋轴**。
\end{mydefinition}

设 \(\boldsymbol{\delta}\) 为任意基点 \(P\) 的位移矢量，\(\boldsymbol{\rho}\) 为旋转的罗德里格斯向量，\(P^*\) 为位于螺旋轴上的一点。则 \(P^*\) 在位移前相对于 \(P\) 的位矢 \(\mathbf{a}^*\) 满足：
\begin{equation}
\mathbf{a}^* = \frac{\boldsymbol{\rho} \times \boldsymbol{\delta} + (\boldsymbol{\rho} \times \boldsymbol{\delta}) \times \boldsymbol{\rho}}{2\boldsymbol{\rho}^2} + \mu \boldsymbol{\rho} \tag{1}
\end{equation}
螺旋平移矢量 \(\boldsymbol{\delta}^*\) 由下式给出：
\begin{equation}
\boldsymbol{\delta}^* = \frac{\boldsymbol{\rho} \cdot \boldsymbol{\delta}}{\boldsymbol{\rho}^2} \boldsymbol{\rho} \tag{2}
\end{equation}

\begin{proof}[公式的证明]
设 \(\mathbf{b}^*\) 为 \(P^*\) 位移后相对 \(P\) 的位矢。则 \(\boldsymbol{\delta}^* = \boldsymbol{\delta} + \mathbf{b}^* - \mathbf{a}^*\)。利用 (1.4.5)，\(\mathbf{b}^* - \mathbf{a}^* = \boldsymbol{\rho} \times (\mathbf{a}^* + \mathbf{b}^*)\)。
若要 \(\boldsymbol{\delta}^*\) 平行于 \(\boldsymbol{\rho}\)，则 \(\boldsymbol{\rho} \times \boldsymbol{\delta}^* = 0\)。对表达式进行代数运算即可解出 \(\mathbf{a}^*\) 与 \(\boldsymbol{\delta}^*\) 的关系，得证式 (1) 和 (2)。
\end{proof}

\section{角速度矩阵 (Angular Velocity Matrix)}

设 \(\mathbf{a}_i\) 和 \(\mathbf{b}_i\) 分别为固定在相对运动的参考系 \(A\) 和刚体 \(B\) 中的单位正交基，方向余弦矩阵 \(C\) 及其元素 \(C_{ij}\) 为时间的函数。
\begin{mydefinition}{角速度矩阵}
定义 \(C\) 的时间导数为 \(\dot{C}\)，则角速度矩阵 \(\widetilde{\boldsymbol{\omega}}\) 定义为：
\begin{equation}
\widetilde{\boldsymbol{\omega}} \triangleq C^T \dot{C} \tag{2}
\end{equation}
由此可得：
\begin{equation}
\dot{C} = C \widetilde{\boldsymbol{\omega}} \tag{3}
\end{equation}
根据反对称矩阵记号，\(\widetilde{\boldsymbol{\omega}}\) 可写为：
\begin{equation}
\widetilde{\boldsymbol{\omega}} = \begin{bmatrix} 0 & -\omega_3 & \omega_2 \\ \omega_3 & 0 & -\omega_1 \\ -\omega_2 & \omega_1 & 0 \end{bmatrix} \tag{4}
\end{equation}
\end{mydefinition}

其中 \(\omega_1, \omega_2, \omega_3\) 由下式给出：
\begin{equation}
\omega_i = \eta_{igh} \dot{C}_{jg} C_{jh} \quad (i=1,2,3) \tag{9}
\end{equation}
其中 \(\eta_{ijk} = \frac{1}{2} \epsilon_{ijk} (\epsilon_{ijk} + 1)\)。反过来，方向余弦矩阵元素的时间导数满足**泊松运动学方程**：
\begin{equation}
\dot{C}_{ij} = C_{ig} \omega_h \epsilon_{ghj} \quad (i,j=1,2,3) \tag{10}
\end{equation}

\begin{proof}[反对称性的证明]
由于 \(C C^T = U\)，对其求导得 \(\dot{C} C^T + C \dot{C}^T = 0\)。左乘 \(C^T\) 得 \(C^T \dot{C} + \dot{C}^T C = 0\)。因为 \(\widetilde{\boldsymbol{\omega}} = C^T \dot{C}\)，所以 \(\widetilde{\boldsymbol{\omega}}^T = -\widetilde{\boldsymbol{\omega}}\)，证毕。
\end{proof}

\section{角速度矢量 (Angular Velocity Vector)}

\begin{mydefinition}{角速度矢量}
矢量 \(\boldsymbol{\omega}\) 定义为：
\begin{equation}
\boldsymbol{\omega} \triangleq \omega_1 \mathbf{b}_1 + \omega_2 \mathbf{b}_2 + \omega_3 \mathbf{b}_3 \tag{1}
\end{equation}
\(\boldsymbol{\omega}\) 称为 \(B\) 在 \(A\) 中的角速度。若用 \({}^A \boldsymbol{\omega}^B\) 表示，则有 \({}^B \boldsymbol{\omega}^A = -{}^A \boldsymbol{\omega}^B\)。
\end{mydefinition}

若用 \(\dot{\mathbf{b}}_i\) 表示 \(\mathbf{b}_i\) 在 \(A\) 中的时间导数，则角速度可以写为：
\begin{equation}
\boldsymbol{\omega} = \mathbf{b}_1 \dot{\mathbf{b}}_2 \cdot \mathbf{b}_3 + \mathbf{b}_2 \dot{\mathbf{b}}_3 \cdot \mathbf{b}_1 + \mathbf{b}_3 \dot{\mathbf{b}}_1 \cdot \mathbf{b}_2 \tag{4}
\end{equation}

\begin{myTheorem}{矢量导数定理}
任意矢量 \(\mathbf{v}\) 在参考系 \(A\) 和 \(B\) 中的时间导数遵循关系：
\begin{equation}
\frac{{}^A d\mathbf{v}}{dt} = \frac{{}^B d\mathbf{v}}{dt} + {}^A \boldsymbol{\omega}^B \times \mathbf{v} \tag{8}
\end{equation}
\end{myTheorem}

\begin{proof}[定理证明]
利用方向余弦矩阵，\(\frac{{}^A d\mathbf{v}}{dt}\) 可写为 \(\dot{^A v} [\mathbf{a}_1\ \mathbf{a}_2\ \mathbf{a}_3]^T\)。利用 \(\dot{^A v} = {^B v} C^T + {^B v} \dot{C}^T\) 展开，代入 \(C\) 的关系，即可证明 (8)。
\end{proof}

\section{角速度分量 (Angular Velocity Components)}

虽然角速度通常分解为固定在 \(B\) 中的矢量，但有时也需要分解到其他参考系。
\begin{myTheorem}{角速度分量的几何解释}
分量 \(\omega_1, \omega_2, \omega_3\) 可以解释为三个特定角度导数在空间中的平均值。即若 \(\alpha\) 是固定在 \(A\) 中的任意单位向量，\(\theta_i\) 是投影相关的角度，则：
\begin{equation}
\omega_i = \overline{\theta}_i = \frac{1}{4\pi} \int \dot{\theta}_i d\sigma \quad (i=1,2,3) \tag{2}
\end{equation}
\end{myTheorem}

\begin{example}[进动与自旋的角速度]
如图 1.12.3，圆柱形航天器 \(B\) 的运动是进动（半顶角 \(\theta\)）与自旋（角度 \(\psi\)）的结合。利用矩阵乘法求得方向余弦矩阵后，由 (1.10.5) 计算角速度分量，可得到其角速度为：
\begin{equation}
\boldsymbol{\omega} = (\dot{\psi} + \dot{\phi} c\theta)\mathbf{b}_1 - \dot{\phi} s\theta c\psi \mathbf{b}_2 + \dot{\phi} s\theta s\psi \mathbf{b}_3 \tag{20}
\end{equation}
将其改写为 \(\boldsymbol{\omega} = \dot{\psi} \mathbf{b}_1 + \dot{\phi} \mathbf{a}_1\)，则两项分别代表了二次简单旋转的角速度，从而物理意义更加清晰。
\end{example}

\section{角速度与欧拉参数 (Angular Velocity and Euler Parameters)}

\begin{myTheorem}{欧拉参数的运动学方程}
设 \(\boldsymbol{\epsilon}\) 和 \(\epsilon_4\) 为欧拉参数，则刚体角速度 \(\boldsymbol{\omega}\) 可以由它们表示：
\begin{equation}
\boldsymbol{\omega} = 2\left( \epsilon_4 \frac{{}^B d\boldsymbol{\epsilon}}{dt} - \dot{\epsilon}_4 \boldsymbol{\epsilon} - \boldsymbol{\epsilon} \times \frac{{}^B d\boldsymbol{\epsilon}}{dt} \right) \tag{1}
\end{equation}
反之，若已知 \(\boldsymbol{\omega}\) 随时间的变化，则可以通过求解以下微分方程得到欧拉参数：
\begin{equation}
\frac{{}^B d\boldsymbol{\epsilon}}{dt} = \frac{1}{2} (\epsilon_4 \boldsymbol{\omega} + \boldsymbol{\epsilon} \times \boldsymbol{\omega}) \tag{2}
\end{equation}
\begin{equation}
\dot{\epsilon}_4 = -\frac{1}{2} \boldsymbol{\omega} \cdot \boldsymbol{\epsilon} \tag{3}
\end{equation}
\end{myTheorem}

利用矩阵形式可以更简洁地表达。定义 \(\boldsymbol{\omega} \triangleq [\omega_1\ \omega_2\ \omega_3\ 0]\)，\(\boldsymbol{\epsilon} \triangleq [\epsilon_1\ \epsilon_2\ \epsilon_3\ \epsilon_4]\)。构造矩阵 \(E\)，则逆关系和正关系分别为：
\begin{equation}
\boldsymbol{\omega} = 2\dot{\boldsymbol{\epsilon}} E \tag{7}
\end{equation}
\begin{equation}
\dot{\boldsymbol{\epsilon}} = \frac{1}{2} \boldsymbol{\omega} E^T \tag{8}
\end{equation}

\section{角速度与罗德里格斯参数 (Angular Velocity and Rodrigues Parameters)}

\begin{myTheorem}{罗德里格斯参数的运动学方程}
设 \(\boldsymbol{\rho}\) 为罗德里格斯向量，则角速度的逆关系为：
\begin{equation}
\boldsymbol{\omega} = \frac{2}{1 + \boldsymbol{\rho}^2} \left( \frac{{}^B d\boldsymbol{\rho}}{dt} - \boldsymbol{\rho} \times \frac{{}^B d\boldsymbol{\rho}}{dt} \right) \tag{1}
\end{equation}
若已知 \(\boldsymbol{\omega}\)，反向求 \(\boldsymbol{\rho}\) 则需要解非线性微分方程：
\begin{equation}
\frac{{}^B d\boldsymbol{\rho}}{dt} = \frac{1}{2} (\boldsymbol{\omega} + \boldsymbol{\rho} \times \boldsymbol{\omega} + \boldsymbol{\rho} \boldsymbol{\rho} \cdot \boldsymbol{\omega}) \tag{2}
\end{equation}
\end{myTheorem}

同样可以写成矩阵形式。定义 \(\widetilde{\boldsymbol{\rho}}\) 为对应的反对称矩阵，则有：
\begin{equation}
\boldsymbol{\omega} = \frac{2\dot{\boldsymbol{\rho}} (U + \widetilde{\boldsymbol{\rho}})}{1 + \boldsymbol{\rho} \boldsymbol{\rho}^T} \tag{6}
\end{equation}
\begin{equation}
\dot{\boldsymbol{\rho}} = \frac{1}{2} \boldsymbol{\omega} (U - \widetilde{\boldsymbol{\rho}} + \boldsymbol{\rho}^T \boldsymbol{\rho}) \tag{7}
\end{equation}

\begin{example}[轴对称航天器的自旋]
考虑轴对称航天器 \(B\) 的自旋问题。通过角动量定理可得到欧拉动力学方程。将动力学方程和运动学方程（式 7）进行数值积分，可以求得罗德里格斯参数随时间的变化。
对于不同的无量纲参数，数值积分的结果在表 1.14.1 中给出。值得注意的是，当 \(\epsilon_4\) 经过零点时，\(\rho_i\) 会趋于无穷大，这凸显了罗德里格斯参数在处理大角度旋转时的奇异性，此时应切换到欧拉参数（如表 1.14.2）。
\end{example}

\section{角速度的间接确定 (Indirect Determination of Angular Velocity)}

当无法直接观测 \(\mathbf{b}_i\) 的导数时，可以通过观测两个非平行向量 \(\mathbf{p}\) 和 \(\mathbf{q}\) 在两个参考系中的变化来求得角速度。
\begin{myformula}[title={间接求解角速度}]
\boldsymbol{\omega} = \frac{ \left( \frac{{}^A d\mathbf{p}}{dt} - \frac{{}^B d\mathbf{p}}{dt} \right) \times \left( \frac{{}^A d\mathbf{q}}{dt} - \frac{{}^B d\mathbf{q}}{dt} \right) }{ \left( \frac{{}^A d\mathbf{p}}{dt} - \frac{{}^B d\mathbf{p}}{dt} \right) \cdot \mathbf{q} }
\end{myformula}
\begin{proof}[公式推导]
基于矢量导数公式 (1.11.8)，\(\frac{{}^A d\mathbf{p}}{dt} - \frac{{}^B d\mathbf{p}}{dt} = \boldsymbol{\omega} \times \mathbf{p}\)，同理对 \(\mathbf{q}\) 也成立。对二者进行叉积，结合矢量三重积展开并化简，即可解出 \(\boldsymbol{\omega}\)。
\end{proof}

\section{辅助参考系 (Auxiliary Reference Frames)}

\begin{myTheorem}{角速度加法定理}
刚体 \(B\) 在参考系 \(A\) 中的角速度可以借助 \(n\) 个辅助参考系 \(A_1, \dots, A_n\) 表示为：
\begin{equation}
{}^A \boldsymbol{\omega}^B = {}^A \boldsymbol{\omega}^{A_1} + {}^{A_1} \boldsymbol{\omega}^{A_2} + \dots + {}^{A_{n-1}} \boldsymbol{\omega}^{A_n} + {}^{A_n} \boldsymbol{\omega}^B \tag{1}
\end{equation}
\end{myTheorem}

对于任意固定于 \(B\) 的矢量 \(\mathbf{c}\)，利用矢量导数定理即可证明此定理。当辅助参考系的运动均为简单旋转时，此定理显得尤为有用。

\section{角速度与方向角 (Angular Velocity and Orientation Angles)}

设三个方向角为 \(\theta_1, \theta_2, \theta_3\)，角速度分量 \(\omega_i = {}^A \boldsymbol{\omega}^B \cdot \mathbf{b}_i\) 与方向角的导数之间由矩阵 \(M\) 关联：
\begin{equation}
[\omega_1 \ \omega_2 \ \omega_3] = [\dot{\theta}_1 \ \dot{\theta}_2 \ \dot{\theta}_3] M \tag{1}
\end{equation}
反之：
\begin{equation}
[\dot{\theta}_1 \ \dot{\theta}_2 \ \dot{\theta}_3] = [\omega_1 \ \omega_2 \ \omega_3] M^{-1} \tag{2}
\end{equation}
针对不同的角度系统，\(M\) 矩阵的具体形式各不相同。
\begin{itemize}
    \item \textbf{空间三角度}：\(M = \begin{bmatrix} 1 & 0 & 0 \\ 0 & c_1 & -s_1 \\ -s_2 & s_1 c_2 & c_1 c_2 \end{bmatrix}\)。奇异性出现在 \(c_2 = 0\)。
    \item \textbf{体三角度}：\(M = \begin{bmatrix} c_2 c_3 & -c_2 s_3 & s_2 \\ s_3 & c_3 & 0 \\ 0 & 0 & 1 \end{bmatrix}\)。奇异性出现在 \(c_2 = 0\)。
    \item \textbf{空间二角度}：\(M = \begin{bmatrix} 1 & 0 & 0 \\ 0 & c_1 & -s_1 \\ c_2 & s_1 s_2 & c_1 s_2 \end{bmatrix}\)。奇异性出现在 \(s_2 = 0\)。
    \item \textbf{体二角度}：\(M = \begin{bmatrix} c_2 & s_2 s_3 & s_2 c_3 \\ 0 & c_3 & -s_3 \\ 1 & 0 & 0 \end{bmatrix}\)。奇异性出现在 \(s_2 = 0\)。
\end{itemize}

\section{慢速小旋转运动 (Slow, Small Rotational Motions)}

当旋转运动不仅幅度小，而且发生得足够慢（即忽略 \(\theta\) 和 \(\dot{\theta}\) 的高阶项）时，动力学方程可以极大地简化。
\begin{align}
\boldsymbol{\omega} &= 2 \frac{{}^B d\boldsymbol{\epsilon}}{dt} \tag{1} \\
\epsilon_4 &= 1 \tag{2} \\
\boldsymbol{\omega} &= 2 \frac{{}^B d\boldsymbol{\rho}}{dt} \tag{3} \\
[\omega_1 \ \omega_2 \ \omega_3] &= [\dot{\theta}_1 \ \dot{\theta}_2 \ \dot{\theta}_3] \tag{4}
\end{align}

\begin{example}[慢小运动的动力学预备推导]
如图 1.18.1，刚体 \(B\) 通过弹性支撑安装在航天器 \(A\) 上。在 \(N\) 系中，\(A\) 的角速度为常量 \(\boldsymbol{\Omega} = \omega_1 \mathbf{a}_1 + \omega_2 \mathbf{a}_2 + \omega_3 \mathbf{a}_3\)。若 \(B\) 相对 \(A\) 为慢小转动，且用体三角度 \(\theta_1, \theta_2, \theta_3\) 描述姿态。为准备建立动力学方程，需计算 \(B\) 相对 \(B^*\) 的角动量 \(\mathbf{H}\) 在 \(N\) 中的时间导数。
利用关系 \(\frac{{}^N d\mathbf{H}}{dt} = \frac{{}^B d\mathbf{H}}{dt} + {}^N \boldsymbol{\omega}^B \times \mathbf{H}\)，并将所有非线性项舍弃（如 \(\theta_i \dot{\theta}_j\) 和 \(\dot{\theta}_i \dot{\theta}_j\) 等），最终可以得到如式 (18) 所示的线性化角动量方程（此处省略了大量冗长展开项）。
\end{example}

\section{瞬时轴 (Instantaneous Axis)}

\begin{mydefinition}{瞬时轴}
当刚体 \(B\) 在参考系 \(A\) 中的角速度 \(\boldsymbol{\omega} \neq 0\) 时，\(B\) 中存在无穷多个点，它们在 \(A\) 中的速度平行于 \(\boldsymbol{\omega}\) 或等于零。这些点构成一条平行于 \(\boldsymbol{\omega}\) 的直线，称为**瞬时轴**。瞬时轴上所有点具有相同的速度 \(\mathbf{v}^*\)，且 \(|\mathbf{v}^*|\) 小于 \(B\) 中任何不在瞬时轴上的点的速度。
\end{mydefinition}

设 \(\mathbf{v}^Q\) 为任意选定的基点 \(Q\) 的速度，\(\boldsymbol{\omega}\) 为角速度，则瞬时轴上一点 \(P^*\) 相对于 \(Q\) 的位矢为：
\begin{equation}
\mathbf{r}^* = \frac{\boldsymbol{\omega} \times \mathbf{v}^Q}{\boldsymbol{\omega}^2} + \mu^* \boldsymbol{\omega} \tag{1}
\end{equation}
其中 \(\mu^*\) 与 \(P^*\) 的选择有关。而 \(\mathbf{v}^*\) 为：
\begin{equation}
\mathbf{v}^* = \frac{\boldsymbol{\omega} \cdot \mathbf{v}^Q}{\boldsymbol{\omega}^2} \boldsymbol{\omega} \tag{2}
\end{equation}

\begin{example}[圆柱卫星的瞬时轴]
如图 1.19.2，做圆周运动的圆柱形卫星，其角速度为 \(\boldsymbol{\omega} = \Omega (\mathbf{c}_1 + \mathbf{c}_2)\)，质心速度为 \(\mathbf{v}^{B^*} = R\Omega \mathbf{c}_1\)。代入公式 (1)，可知瞬时轴穿过质心 \(B^*\) 与轨道中心 \(A^*\) 的中点，并分别与 \(\mathbf{c}_1\) 和 \(\mathbf{c}_2\) 成 \(45^\circ\) 夹角。
\end{example}

\section{角加速度 (Angular Acceleration)}

\begin{mydefinition}{角加速度}
刚体 \(B\) 在参考系 \(A\) 中的角加速度 \(\boldsymbol{\alpha}\) 定义为角速度 \(\boldsymbol{\omega}\) 在 \(A\) 中的一阶时间导数：
\begin{equation}
\boldsymbol{\alpha} \triangleq \frac{{}^A d\boldsymbol{\omega}}{dt} \tag{1}
\end{equation}
\end{mydefinition}

若将角速度和角加速度分解到固定的参考系 \(C\)（其自身具有角速度 \(\boldsymbol{\Omega}\)）中，即 \(\boldsymbol{\omega} = \sum {}^C\omega_i \mathbf{c}_i\)，\(\boldsymbol{\alpha} = \sum {}^C\alpha_i \mathbf{c}_i\)，则有：
\begin{equation}
{}^C \alpha_i = {}^C \dot{\omega}_i + \boldsymbol{\Omega} \times \boldsymbol{\omega} \cdot \mathbf{c}_i \quad (i=1,2,3) \tag{4}
\end{equation}
这说明角加速度的分量分量不仅取决于对应角速度分量的导数，还与参考系 \(C\) 自身的旋转有关。

\section{偏角速度和偏速度 (Partial Angular Velocities and Partial Velocities)}

\begin{mydefinition}{广义速度与偏速度}
在动力学分析中，经常引入 \(n\) 个**广义速度** \(u_1, \dots, u_n\)，它们是广义坐标导数的线性组合。此时，刚体 \(B\) 的角速度 \(\boldsymbol{\omega}\) 和质点 \(P\) 的速度 \(\mathbf{v}\) 可以唯一地表示为：
\begin{align}
\boldsymbol{\omega} &= \sum_{r=1}^{n} \boldsymbol{\omega}_r u_r + \boldsymbol{\omega}_t \tag{2} \\
\mathbf{v} &= \sum_{r=1}^{n} \mathbf{v}_r u_r + \mathbf{v}_t \tag{3}
\end{align}
其中 \(\boldsymbol{\omega}_r\) 称为 \(B\) 在 \(A\) 中的第 \(r\) 个**偏角速度**，\(\mathbf{v}_r\) 称为 \(P\) 在 \(A\) 中的第 \(r\) 个**偏速度**。\(\boldsymbol{\omega}_t\) 和 \(\mathbf{v}_t\) 仅依赖于坐标和时间，不包含广义速度 \(u_r\)。
\end{mydefinition}

\begin{proof}[偏速度矩阵形式的证明]
利用方向余弦矩阵 \(C\) 的导数关系 \(\dot{C} = C \widetilde{\boldsymbol{\omega}}\)，将 \(u_r\) 作为独立变量代入，按 \(u_r\) 提取系数，即得到 \(\boldsymbol{\omega}_r\) 的形式。
\end{proof}

\begin{example}[圆周运动刚体的偏速度]
图 1.21.1 中，刚体 \(B\) 绕圆轨道以恒定速度 \(V\) 运动，同时自身旋转。若定义广义速度 \(u_r = \boldsymbol{\omega} \cdot \mathbf{b}_r\)，则 \(\boldsymbol{\omega} = u_1 \mathbf{b}_1 + u_2 \mathbf{b}_2 + u_3 \mathbf{b}_3\)。对应的偏角速度分别为 \(\boldsymbol{\omega}_1 = \mathbf{b}_1\)，\(\boldsymbol{\omega}_2 = \mathbf{b}_2\)，\(\boldsymbol{\omega}_3 = \mathbf{b}_3\)。若换成空间三角度作为广义坐标，虽然推导思路相同，但表达式的形式会变得极其复杂（如式 22-25 所示）。因此，合理选择广义速度至关重要。
\end{example}